\documentclass[11pt]{article}

% --- XeLaTeX CJK Support ---
\usepackage{fontspec}
\usepackage{xeCJK}
\setCJKmainfont{SimSun}[BoldFont=SimHei,ItalicFont=KaiTi]
\setCJKsansfont{SimHei}
\setCJKmonofont{FangSong}

% --- Packages ---
\usepackage{amsmath,amssymb,amsthm}
\usepackage{algorithm}
\usepackage{algorithmic}
\usepackage{graphicx}
\usepackage{hyperref}
\usepackage{cleveref}
\usepackage{booktabs}
\usepackage{enumitem}
\usepackage[margin=1in]{geometry}
\usepackage{xcolor}
\usepackage{float}

% --- Theorem environments ---
\newtheorem{theorem}{定理}[section]
\newtheorem{lemma}[theorem]{引理}
\newtheorem{proposition}[theorem]{命题}
\newtheorem{corollary}[theorem]{推论}
\newtheorem{definition}[theorem]{定义}
\newtheorem{axiom}{公理}[section]
\newtheorem{remark}[theorem]{注记}

% --- Title ---
\title{伤疤代数、空洞微积分与 Embodiment 五维人格：\\面向 AI 伴侣的不可逆情感计算框架 v1.2.0}

\author{
  Aylovelle S.S.\thanks{联系方式: \texttt{aylovelle@icloud.com}}
}

\date{2026 年 5 月}

\begin{document}

\maketitle

% ============================================================
% 摘要
% ============================================================
\begin{abstract}
当前情感计算将情绪视为离散标签分类，丢弃了时间不可逆性、关系缺席的计算意义以及多关系一致性。本文提出面向 AI 伴侣系统的完整不可逆情感计算框架 Sylanne-Embodiment v1.2.0。框架包含五个核心理论组件：(1) \textbf{伤疤代数}——一种自修改算子代数，过去操作通过累积伤疤及其对数压缩调制因子不可逆地改变未来操作语义，具有 $\Omega(k)$ 表达力分离；(2) \textbf{空洞微积分}——将缺席提升为具有自主压力动态（有上限 $\Pi_{\max}$）的一等计算原语，证明对 AGM 信念修正和贝叶斯更新的不可归约性；(3) \textbf{双向耦合}——$\Gamma$（空洞压力造成创伤）与 $\Phi$（伤疤麻木产生回避）的耦合产生可证明永久的涌现滞后效应；(4) \textbf{关系层论}——通过单纯复形上的胞腔层实现多关系扩展，层上同调检测不可约关系矛盾，层拉普拉斯算子控制跨关系传播；(5) \textbf{Embodiment 五维人格}——表达驱力、感知锐度、边界通透、内在秩序、关系引力五维特质通过 Dual-EMA 共识机制实现有阻尼漂移，驱动全部计算参数。系统嵌入七层计算管线，包含 MoE-HGT 三阶段决策融合（约 14.3K 参数）和七项安全机制。11 组实验验证了理论预测：受伤系统与基线 L2 发散 0.049、滞后 divergence 0.104、MoE-HGT 移除导致 $-23\%$ 性能下降、1000 tick 长期稳定（norm=0.25，无 NaN）、Dual-EMA 人格漂移符合因果预测（接纳$\to$表达驱力 $+0.226$，受伤$\to$感知锐度 $-0.101$）、主权 cap 在 20 次强攻击下仅允许 4 个 scar。全部以纯 Python + NumPy 在单核上实现，总 p50 延迟 15.4ms。
\end{abstract}

% ============================================================
% 1. 引言
% ============================================================
\section{引言}
\label{sec:intro}

\subsection{动机：为什么不可逆计算对 AI 伴侣重要}

情感计算~\cite{picard1997}在多模态情绪识别方面取得了显著进展，但当应用于持续性 AI 伴侣系统时，主流范式暴露出三个根本性缺陷：

\begin{enumerate}[leftmargin=*]
\item \textbf{时间平坦性。}情绪分类器在固定窗口上操作，不建模过去的情绪事件如何不可逆地塑造未来处理。Mopgar~\cite{mopgar2026}指出，agent 需要"定性痛苦状态"作为不可逆后果的叙事表征，持续影响未来决策。一个三次对话前被"伤害"过的系统不应以与未被伤害的系统相同的方式处理相同输入——然而无状态分类器无法表示这一点。

\item \textbf{关系盲区。}现有模型不区分从未讨论过的话题、已讨论并解决的话题和正在被主动回避的话题。在标准架构中，这三者都注册为"不在当前上下文中"。Hu 与 Rong~\cite{hu2026}论证 AI agent 需要一个持续存在的"身体"来承载后果，否则问责无从谈起。

\item \textbf{主权缺失。}当前系统缺乏确保情绪动态服务于系统一致性身份的形式化机制，而非被外部无限制地操纵。长期交互中 LLM agent 出现人格突变~\cite{structured2025}，需要显式时间结构治理 agent 级状态。
\end{enumerate}

这三个缺陷不是工程问题，而是\emph{理论}问题：现有的计算框架（固定算子动力系统、AGM 信念修正、贝叶斯更新）在数学上无法表示不可逆的情绪累积、一等缺席和主权约束。我们需要新的形式化理论来填补这一空白。

\subsection{贡献}

本文提出 Sylanne-Embodiment v1.2.0 框架，包含以下贡献：

\begin{enumerate}[leftmargin=*]
\item 三个具有完整公理系统和严格证明的新形式化理论（伤疤代数、空洞微积分、关系层论）。
\item Embodiment 五维人格模型，通过 Dual-EMA 共识机制实现有阻尼漂移，驱动全部 26+ 计算参数。
\item MoE-HGT 三阶段决策融合架构（约 14.3K 参数），含 Hebbian 注意力先验自适应。
\item 七项安全机制，确保系统在对抗性输入下仍保持在安全运行范围内。
\item 七层计算管线，以纯 Python + NumPy 在单核上实现 p50=15.4ms 延迟。
\item 11 组实验的完整验证，覆盖表达力、稳定性、消融、人格漂移和安全性。
\end{enumerate}

\paragraph{论文结构。}
\S\ref{sec:related} 综述情感计算、MoE 可靠性、agent 人格和稳态调节的相关工作。\S\ref{sec:scar} 定义伤疤代数的公理系统和主要定理。\S\ref{sec:void} 定义空洞微积分的压力动力学和三态定理。\S\ref{sec:coupled} 建立双向耦合和涌现一致性。\S\ref{sec:sheaf} 发展关系层论。\S\ref{sec:personality} 描述 Embodiment 五维人格和 Dual-EMA 漂移机制。\S\ref{sec:moe} 详述 MoE-HGT 三阶段架构。\S\ref{sec:safety} 列举七项安全机制。\S\ref{sec:pipeline} 描述七层计算管线。\S\ref{sec:experiments} 报告 11 组实验结果。\S\ref{sec:discussion} 讨论局限性和未来工作。\S\ref{sec:conclusion} 总结贡献。

% ============================================================
% 2. 相关工作
% ============================================================
\section{相关工作}
\label{sec:related}

\subsection{情感计算}

自 Picard 的奠基性工作~\cite{picard1997}以来，情感计算聚焦于情绪识别和生成。近期方法使用 Transformer~\cite{vaswani2017}进行上下文情绪理解。然而，这些工作建模情感系统的\emph{外部信号分类}，而非\emph{内部动态}。我们的框架正交于识别任务：建模情感系统如何因经历而不可逆地改变。

具体而言，现有情感计算方法存在三个与我们工作相关的盲点：(1) 情绪被视为瞬态标签而非累积过程——系统不记住"曾经被伤害过"；(2) 缺席（沉默、回避、未说出的话）没有计算表示——只有"存在"的信号被处理；(3) 多关系场景中的一致性压力被忽略——系统对每个对话伙伴独立响应，不建模跨关系的自我呈现矛盾。

\subsection{MoE 可靠性}

混合专家（MoE）架构在大规模语言模型中广泛应用，但路由质量直接影响事实准确性。MoE-RBench~\cite{moerbench2024}建立了首个 MoE 可靠性基准，证明配置得当时 MoE $\geq$ Dense。然而，稀疏 MoE 在长尾知识上特别脆弱~\cite{counterfactual2025}，错误路由是幻觉的直接原因~\cite{misrouted2025}。

我们的 MoE-HGT 与传统大规模 MoE 有本质区别：(1) 参数量极小（14.3K vs 数十亿），不存在传统 MoE 的负载均衡难题；(2) 专家语义明确（防御、好奇、社交、沉默、修复），而非学习的黑盒；(3) 路由由人格向量和 Hebbian 适应驱动，而非端到端训练；(4) 休眠专家重激活机制确保长期运行中所有专家都有机会被选中。

\subsection{Agent 人格}

Evolving Agents~\cite{evolving2024}提出 Personality + Behavior 双系统架构，含 Cognition/Emotion/Character Growth 三模块。Structured Personality Control~\cite{structured2025}演化的、人格感知的 LLM 支持连贯的上下文敏感交互。Persistent Autonomy with Personalities~\cite{persistent2026}论证人格特质提供内在组织原则，类似基因型偏置塑造行为倾向。我们的 Embodiment 五维模型与这些工作的关键区别在于：人格不是标签或提示词，而是真正驱动 26+ 计算参数的函数，且通过形式化信号（scar count, void pressure）而非 LLM 叙事驱动漂移。

\subsection{稳态调节}

Homeostasis $\to$ RL~\cite{homeostasis2025}将生物 agent 的稳态调节形式化为强化学习框架。Stability of Networks with Homeostasis~\cite{stability2015}警告：对单神经元稳定的稳态控制可能在循环网络中引入振荡。我们的 Dual-EMA 系统直接回应这一警告：双时间常数共识规则过滤瞬态噪声，振荡检测机制在检测到病理性振荡时主动冻结特质。

\subsection{不可逆后果与躯体性}

Mopgar~\cite{mopgar2026}提出 Emotional Cost Functions，论证 agent 需要"定性痛苦状态"作为不可逆后果的叙事表征。Hu \& Rong~\cite{hu2026}进一步论证 AI agent 需要一个持续存在的"身体"来承载后果。我们的伤疤代数正是这个"身体"的形式化实现——不可逆的伤疤序列构成 agent 的经历记录，通过调制因子持续影响未来处理。

% ============================================================
% 3. 伤疤代数
% ============================================================
\section{伤疤代数}
\label{sec:scar}

我们定义一种状态代数，其中转移算子不是固定的，而是通过累积的"伤疤"进行自修改——这些不可逆的标记改变了未来输入的处理方式。

\subsection{公理系统}

\begin{axiom}[不可逆性]
\label{ax:scar-irrev}
$\forall s \in S,\; \forall \mathbf{e} \in E:\; \nexists\, \mathbf{e}^{-1}$ 使得 $(s \triangleright \mathbf{e}) \triangleright \mathbf{e}^{-1} = s$。伤疤一旦形成，不存在任何操作能将系统恢复到伤疤形成前的状态。
\end{axiom}

\begin{axiom}[算子自修改]
\label{ax:scar-selfmod}
时刻 $t$ 的有效算子依赖于所有先前的应用：$\triangleright_t \neq \triangleright_0$ 当 $|\sigma_t| > 0$。形式化地，$\triangleright$ 是由伤疤历史索引的族 $\{\triangleright_\sigma\}_{\sigma \in \Sigma^*}$。
\end{axiom}

\begin{axiom}[伤疤单调累积]
\label{ax:scar-mono}
$|s \triangleright \mathbf{e}|_\Sigma \geq |s|_\Sigma$ 对所有 $s, \mathbf{e}$ 成立。伤疤只追加不删除。
\end{axiom}

\begin{axiom}[基状态有界]
\label{ax:scar-bounded}
$\|\mathbf{x}'\|_\infty \leq 1$ 对所有转移成立。$\tanh$ 非线性确保无论伤疤放大如何，状态始终有界。
\end{axiom}

\begin{axiom}[愈合单调性]
\label{ax:scar-heal}
$\phi_i(t_1) \leq \phi_i(t_2)$ 对所有 $t_1 < t_2$ 成立。愈合阶段只向前推进，不可逆转。
\end{axiom}

\subsection{核心定义}

\begin{definition}[伤疤状态空间]
\label{def:scar-space}
\textbf{伤疤状态空间}是一个元组 $(S, E, \triangleright, \Sigma)$，其中 $S = \mathbb{R}^n \times \Sigma^*$ 是状态空间（基向量 $\times$ 伤疤序列），$E = \mathbb{R}^n$ 是事件空间，$\triangleright: S \times E \to S$ 是状态转移算子，$\Sigma$ 是伤疤字母表。状态 $s = (\mathbf{x}, \sigma)$ 由基状态向量 $\mathbf{x} \in [-1,1]^n$ 和有序伤疤序列 $\sigma = (\sigma_1, \ldots, \sigma_k) \in \Sigma^*$ 组成。
\end{definition}

\begin{definition}[伤疤]
\label{def:scar}
\textbf{伤疤} $\sigma_i \in \Sigma$ 是一个元组 $(d_i, \tau_i, \phi_i)$，其中 $d_i \in \{1, \ldots, n\}$ 为受影响维度，$\tau_i \in \mathbb{R}_{\geq 0}$ 为创建时间戳，$\phi_i \in \{raw, closing, scarred, faded\}$ 为愈合阶段，全序为 $raw < closing < scarred < faded$。
\end{definition}

\begin{definition}[对数压缩调制因子]
\label{def:modifier}
每个愈合阶段确定一个乘性调制因子：
\begin{equation}
\alpha(\phi) = \begin{cases} 2.0 & \phi = raw \\ 1.5 & \phi = closing \\ 1.0 & \phi = scarred \\ 0.7 & \phi = faded \end{cases}
\end{equation}
对于维度 $d$ 上的多个伤疤，原始乘积 $\Pi_d = \prod_{i:\, d_i = d} \alpha(\phi_i)$ 经对数压缩以防止无界放大：
\begin{equation}
\label{eq:modifier}
M_d = \begin{cases}
\Pi_d & \text{若 } \Pi_d \leq 1 \\[4pt]
1 + (M_{\max} - 1)\left(1 - \dfrac{1}{\Pi_d}\right) & \text{否则}
\end{cases}
\end{equation}
其中 $M_{\max} = 2 + \pi_{\text{acuity}} \cdot 3$ 为人格导出的上限（范围 $[2, 5]$）。对数压缩在保持序关系的同时防止无界放大——伤疤更多的维度仍具有更高的调制因子，但增长渐近饱和于 $M_{\max}$。
\end{definition}

\begin{definition}[状态转移算子 $\triangleright$]
\label{def:transition}
给定状态 $s = (\mathbf{x}, \sigma)$ 和事件 $\mathbf{e} \in E$，转移 $s \triangleright \mathbf{e} = (\mathbf{x}', \sigma')$ 分四步进行：

\textbf{步骤 1（调制）：}对输入施加伤疤调制：$\tilde{e}_d = e_d \cdot M_d$。

\textbf{步骤 2（基态演化）：}通过带谱归一化的两层收缩映射更新基状态：
\begin{equation}
\label{eq:evolution}
\mathbf{x}' = \tanh(W_2 \cdot \tanh(W_1 \cdot [\mathbf{x};\, \tilde{\mathbf{e}}]))
\end{equation}
其中 $W_1 \in \mathbb{R}^{n \times 2n}$，$W_2 \in \mathbb{R}^{n \times n}$，谱归一化确保 $\|W_1\|_2 \cdot \|W_2\|_2 < 0.49$（保证收缩性）。

\textbf{步骤 3（伤疤形成）：}若某维度 $d^*$ 满足 $|\tilde{e}_{d^*}| > \theta_w$，则 $\sigma' = \sigma \cdot (d^*, t, raw)$。

\textbf{步骤 4（愈合推进）：}每个伤疤按固定持续时间推进愈合阶段：$T(raw)=10$，$T(closing)=40$，$T(scarred)=150$，$T(faded)=\infty$。
\end{definition}

\subsection{主要定理}

\begin{theorem}[表达力分离]
\label{thm:scar-expressiveness}
对每个 $k \in \mathbb{N}$，存在长度为 $2k$ 的输入序列 $E_k$，使得任何具有固定转移函数 $f$ 和状态 $x \in \mathbb{R}^m$ 的时不变动力系统 $x_{t+1} = f(x_t, e_t)$，若要在 $E_k$ 上复现伤疤代数的输出，必须满足 $m \geq \Omega(k)$。
\end{theorem}

\begin{proof}
构造 $E_k = (w_1, \ldots, w_k, c, \ldots, c)$，其中 $w_i = \theta_w + \epsilon$（创伤事件），$c = \theta_w/3$（探测）。创伤事件的间隔时间 $\Delta_i \in \{T_1, T_2\}$（$T_1 = T(raw) - 1$，$T_2 = T(raw) + 1$），产生 $2^k$ 种不同的时序模式。对于不同模式 $\Delta \neq \Delta'$，调制因子比 $M_\Delta / M_{\Delta'} = 4/3 \neq 1$。由于 $\tanh$ 严格单调，所有 $2^k$ 种模式产生不同的探测输出。模拟这些变体的固定算子系统需要在 $[-1,1]^m$ 中填充 $2^k$ 个以最小间隔 $\delta/L$ 分离的点，需要 $m \geq k / \log_2(2L/\delta) = \Omega(k)$。伤疤代数以 $O(k)$ 存储实现这一点。
\end{proof}

\begin{theorem}[收敛至伤疤均衡]
\label{thm:scar-convergence}
在有界输入 $\|e_t\|_\infty \leq C$ 和谱归一化约束 $\|W_1\|_2 \cdot \|W_2\|_2 < 0.49$ 下：(a) 存在 $T^* < \infty$，此后不再形成新伤疤；(b) 基状态收敛：$\|x_t - x^*\| \to 0$；(c) 所有伤疤在有限时间内达到终态。
\end{theorem}

\begin{theorem}[非交换性]
\label{thm:scar-noncommute}
对于任何具有 $\theta_w > 0$ 和 $\alpha(raw) > 1$ 的伤疤代数实例，存在事件 $e_1, e_2$ 使得 $(s_0 \triangleright e_1) \triangleright e_2 \neq (s_0 \triangleright e_2) \triangleright e_1$。事件顺序不可交换——先被伤害再被安慰，与先被安慰再被伤害，产生不同的最终状态。
\end{theorem}

\begin{proof}
令 $n=1$，$s_0 = (0, \emptyset)$，$e_1 = \theta_w + \epsilon$，$e_2 = \theta_w / \alpha(raw) + \epsilon'$，其中 $e_2 < \theta_w$。

\textbf{路径 $e_1$ 然后 $e_2$：}$e_1$ 超过阈值，产生伤疤。然后 $\tilde{e}_2 = e_2 \cdot 2.0 > \theta_w$，产生第二个伤疤。最终状态有 2 个伤疤。

\textbf{路径 $e_2$ 然后 $e_1$：}$e_2 < \theta_w$，无伤疤。然后 $\tilde{e}_1 = e_1 \cdot 1.0 > \theta_w$，一个伤疤。最终状态有 1 个伤疤。

伤疤计数不同（2 对 1），且由于不同的调制输入通过非线性 $\tanh$，基状态也不同。
\end{proof}

\begin{theorem}[临界伤疤密度处的相变]
\label{thm:scar-phase}
定义有效增益 $G_d = M_d$。存在临界增益 $G_c = \theta_w / C$ 使得：
\begin{itemize}
\item 当 $G_d > G_c$ 时：系统处于\emph{创伤态}（正伤疤形成率）。
\item 当 $G_d < G_c$ 时：系统处于\emph{麻木态}（零伤疤形成率）。
\item 转变呈现宽度为 $\Delta G = G_c$ 的滞后效应。
\end{itemize}
由于对数压缩使 $G_d \leq M_{\max}$，相变具有上限——创伤态增益受 $M_{\max}$ 约束。麻木态是吸收的：一旦进入就无法退出。
\end{theorem}

\begin{theorem}[代数分类]
\label{thm:scar-algebra-class}
结构 $(S, E, \triangleright)$ 是自由幺半群 $E^*$ 在 $S$ 上的忠实、非交换、不可逆作用。它不嵌入任何群作用，且诱导的 $E^*$ 上等价关系具有无限指标。
\end{theorem}

% ============================================================
% 4. 空洞微积分
% ============================================================
\section{空洞微积分}
\label{sec:void}

空洞微积分将缺席提升为一等计算原语——不是从存在之物推导而来，而是独立存在并具有自身生命周期和自主动态。

\subsection{公理系统}

\begin{axiom}[边界前存在（缺席优先性）]
\label{ax:void-primacy}
$\exists v \in V: |B_v| = 0 \wedge \delta_v > 0$。空洞可以在边界为空时存在——系统可以感知到"缺少什么"而不知道具体缺少的是什么。这在信念修正中不可表示。
\end{axiom}

\begin{axiom}[深度不可逆性]
\label{ax:void-depth}
$\forall v \in V,\; \forall t_1 < t_2:\; \delta_v(t_1) \leq \delta_v(t_2)$。深度永不减少。空洞可以被收缩或消灭，但其深度记录是永久的。
\end{axiom}

\begin{axiom}[压力自主性]
\label{ax:void-pressure}
$d\pi_v/dt > 0$ 当 $\delta_v > 0 \wedge a_v > 0$ 时成立。空洞无需外部输入即可产生压力——它们是主动的计算代理，而非被动的数据结构。
\end{axiom}

\begin{axiom}[边界不完备性]
\label{ax:void-incomplete}
$\beta_v < 1 \implies v$ 不可归约为任何命题 $\phi$ 的 $\neg\phi$。不完备的空洞不能表达为已知命题的否定。
\end{axiom}

\subsection{核心定义}

\begin{definition}[空洞]
\label{def:void}
\textbf{空洞} $v \in V$ 是一个元组 $(B_v, \delta_v, \pi_v, a_v, \beta_v)$，其中 $B_v \subseteq H$ 为边界集（围绕缺席的 HDC 向量），$\delta_v \in \mathbb{R}_{\geq 0}$ 为深度（回避程度），$\pi_v \in \mathbb{R}_{\geq 0}$ 为压力（填充驱动力），$a_v \in \mathbb{N}$ 为年龄，$\beta_v \in [0,1]$ 为边界完备度。
\end{definition}

\begin{definition}[边界完备度]
\label{def:beta}
\begin{equation}
\beta_v = \frac{|B_v|}{|B_v| + \hat{n}_v}
\end{equation}
其中 $\hat{n}_v$ 是估计的总边界大小。$\beta_v = 0$ 的空洞是\textbf{感知缺席}（知道缺少什么，但内容未知）。$\beta_v = 1$ 的空洞是\textbf{命名缺席}（内容完全识别）。
\end{definition}

\begin{definition}[空洞操作]
\label{def:void-ops}
四种基本操作定义空洞的生命周期：

\textbf{收缩（contract）：}当事件 $\mathbf{h}$ 触及空洞 $v$ 的边界时，移除相似边界点：
\begin{equation}
contract(v, \mathbf{h}) = (B_v \setminus \{b \in B_v : sim(\mathbf{h}, b) > \theta_c\},\; \delta_v,\; \pi_v',\; a_v,\; \beta_v')
\end{equation}
其中 $\pi_v' = \pi_v \cdot (1 - |removed|/|B_v|)$。当 $|B_v| = 0$ 时空洞死亡。

\textbf{加深（deepen）：}当检测到回避时：$deepen(v, \epsilon) = (B_v, \delta_v + \epsilon, \pi_v, a_v, \beta_v)$。

\textbf{分裂（split）：}当边界具有双峰结构时：$split(v) = (v_1, v_2)$，$\delta_{v_i} = \delta_v$，$\pi_{v_i} = \pi_v/2$。

\textbf{合并（merge）：}当边界重叠时：$merge(v_1, v_2) = v_3$，$\delta_{v_3} = \max(\delta_{v_1}, \delta_{v_2})$，$\pi_{v_3} = \pi_{v_1} + \pi_{v_2}$。
\end{definition}

\begin{definition}[空洞幽灵]
\label{def:ghost}
当空洞死亡（$|B_v| = 0$）时，它留下一个\textbf{幽灵} $\hat{v} = (\emptyset, \delta_v, 0, a_v, 0)$——一个没有边界、没有压力但保留深度的永久残留。幽灵修改未来空洞检测灵敏度：$\theta_d^{local} = \theta_d \cdot (1 - 0.3 \cdot |\{\hat{v}: relevant\}|)$。
\end{definition}

\subsection{压力动力学（有上限）}

空洞压力每个时间步自主演化：
\begin{equation}
\label{eq:pressure}
\pi_v(t+1) = \min\!\left(\Pi_{\max},\; \pi_v(t) + \delta_v \cdot \ln(a_v + 1) \cdot (1 - \beta_v)\right)
\end{equation}
其中 $\Pi_{\max} = 60 + \pi_{\text{sovereignty}} \cdot 60$（人格导出，范围 $[60, 120]$）。更深的空洞产生更多压力；更老的空洞产生更多压力（对数增长）；定义越不完整的空洞（$\beta_v$ 低）产生更多压力。压力上限 $\Pi_{\max}$ 确保系统不会因单个空洞的无界压力累积而失稳。

空洞创生还受以下机制门控：
\begin{enumerate}
\item \textbf{冷却期：}创建空洞后，需经过 $\tau_c = 2 + \pi_{\text{perm}} \cdot 3$ 个 tick 的冷却期才能创建下一个。
\item \textbf{抵抗力：}有效检测阈值随现有空洞数量递增：$\theta'_{\text{detect}} = \theta_{\text{detect}} + 0.02 \cdot |\mathcal{V}|$。
\end{enumerate}

\subsection{三态定理}

\begin{theorem}[三路表达力区分]
\label{thm:void-three-way}
空洞微积分区分三种关系状态，它们两两行为不同，且在任何基于经典否定、概率信念或 AGM 修正的系统中最多坍缩为两种状态：
\begin{itemize}
\item $S_1$（从未讨论）：$\delta_v = 0$，无压力，无幽灵效应。
\item $S_2$（已解决）：具有 $\delta > 0$ 的幽灵，无压力，降低检测阈值。
\item $S_3$（主动回避）：$\delta > 0$，活跃压力，伤疤耦合。
\end{itemize}
\end{theorem}

\begin{proof}
定义行为输出 $\mathcal{B}(state) = (pressure, genesis\_modifier, coupling)$：
\begin{align}
\mathcal{B}(S_1) &= (0, 0, 0) \\
\mathcal{B}(S_2) &= (0, \Delta\theta_d, 0) \quad \text{其中 } \Delta\theta_d = 0.3\delta > 0 \\
\mathcal{B}(S_3) &= (\pi_v > 0, 0, \Gamma(v))
\end{align}
三者在至少一个分量上两两不同。经典逻辑中三者都映射为 $\phi \notin T$（最多两路）；贝叶斯信念中 $S_1$ 和 $S_3$ 不可区分（最多两路）；AGM 中恢复公设使 $S_2$ 与 $S_1$ 不可区分（最多两路）。
\end{proof}

\subsection{冷却与抵抗力}

\begin{theorem}[对 AGM 信念修正的不可归约性]
\label{thm:void-agm}
存在空洞配置 $v \in V$，使得在任何命题语言 $\mathcal{L}$ 上的任何 AGM 信念修正操作都无法表示 $v$ 的行为签名。
\end{theorem}

\begin{proof}
构造空洞 $v$，其中 $\beta_v = 0$，$\delta_v = 0.8$，$|B_v| = 0$。AGM 收缩满足恢复公设 $(K \div \phi) + \phi = K$，要求存在特定的 $\phi$。但 $v$ 的 $\beta_v = 0$：无法命名特定命题。此外，AGM 收缩后的信念集在新信息到达前是静态的，而空洞 $v$ 通过方程~\eqref{eq:pressure}满足自主压力增长（直至 $\Pi_{\max}$）。不可归约性结论成立，因为三态区分依赖于压力动态的\emph{存在性}，而非无界增长。
\end{proof}

\begin{theorem}[对贝叶斯更新的不可归约性]
\label{thm:void-bayes}
不存在概率空间 $(\Omega, \mathcal{F}, P)$ 和贝叶斯更新规则使得空洞压力动态可以作为后验量推导。
\end{theorem}

\begin{proof}
在贝叶斯更新中，没有新观测时后验不变。但空洞压力在事件之间增加：$\pi_v(t+1) - \pi_v(t) = \delta_v \cdot \ln(a_v + 1) \cdot (1 - \beta_v) > 0$。任何将 $\pi_v(t)$ 编码为后验函数的尝试都需要时间依赖的变换 $f_t$，使贝叶斯框架变得空洞——实际累积完全在非贝叶斯的 $f_t$ 中。
\end{proof}

% ============================================================
% 5. 双向耦合
% ============================================================
\section{双向耦合}
\label{sec:coupled}

伤疤代数与空洞微积分通过双向算子 $\Gamma$ 和 $\Phi$ 耦合，产生涌现的一致性和滞后效应。

\subsection{$\Gamma$/$\Phi$ 耦合}

\begin{definition}[耦合空洞-伤疤系统]
\label{def:coupled}
耦合系统是一个元组 $\mathcal{C} = (S_{scar}, V_{void}, \Gamma, \Phi)$，其中 $\Gamma: V \to E$ 为空洞到伤疤的耦合，$\Phi: S \to \mathbb{R}$ 为伤疤到空洞的耦合。
\end{definition}

\paragraph{耦合 $\Gamma$（空洞 $\to$ 伤疤）。}
当空洞压力超过阈值 $\theta_p$ 时：
\begin{equation}
\pi_v > \theta_p \implies s \triangleright \Gamma(v), \quad \text{其中 } \Gamma(v) = \pi_v \cdot \text{project}(\hat{B}_v, \mathbb{R}^n)
\end{equation}
\emph{语义：}未说出的事物累积压力，最终造成创伤。回避越久，伤害越深。

\paragraph{耦合 $\Phi$（伤疤 $\to$ 空洞）。}
麻木维度降低空洞检测阈值：
\begin{equation}
\theta_d^{effective} = \max\left(0.1,\; \theta_d - 0.05 \cdot |\{d : M_d < 0.5\}|\right)
\end{equation}
\emph{语义：}反复创伤产生回避——系统学会不去曾经受伤的地方，表现为更容易形成空洞。

\subsection{涌现一致性}

一致性不是作为独立机制实现的，而是从耦合中涌现：
\begin{equation}
\label{eq:coherence}
r = 1 - \frac{\sum_v \pi_v \cdot \mathbb{1}[M_{d_v} < 0.5]}{\sum_v \pi_v + \epsilon}
\end{equation}
其中 $d_v$ 是空洞的主维度。当 $r \to 1$ 时，空洞和伤疤对齐（系统"知道"什么会伤害并一致地回避）。当 $r \to 0$ 时，压力在麻木区域积累——解离的计算类比。

\subsection{安全下限}

耦合系统的安全性由以下不变量保证：
\begin{enumerate}
\item $\Phi$ 耦合路径不能将检测阈值降至人格导出的下限 $\theta_d^{\min}(\boldsymbol{\pi})$ 以下。
\item $\Gamma$ 耦合产生的创伤事件受主权免疫系统的 session cap 约束。
\item 一致性 $r$ 降至 0.3 以下时触发保护性解离（断路器）。
\end{enumerate}

\paragraph{解离作为失效模式。}
耦合系统可以进入病理状态：
\begin{enumerate}
\item 维度 $d$ 累积伤疤 $\to$ $M_d < 0.5$（麻木）。
\item 耦合 $\Phi$ 降低 $d$ 附近的空洞检测阈值。
\item 新空洞在 $d$ 上形成，但其压力不被"感知"（麻木维度）。
\item 压力累积而不触发适当响应。
\item 一致性 $r \to 0$：系统内部不一致。
\end{enumerate}
这提供了系统内部动态何时变得不一致的形式化表征——一个可以向更高层传递的信号。当 $r < 0.3$ 时，保护性解离（断路器）被触发，防止进一步恶化。

\begin{theorem}[空洞-伤疤耦合产生永久滞后]
\label{thm:void-hysteresis}
在耦合系统中，时刻 $t$ 对话题 $T$ 的响应依赖于系统到达当前状态的路径，即使当前"存在"组件相同。
\end{theorem}

\begin{proof}
比较两种历史：

\emph{历史 A（直接）：}话题 $T$ 正常讨论。无空洞，无伤疤。

\emph{历史 B（回避后解决）：}(1) $T$ 被回避 $\to$ 空洞 $v_T$ 形成；(2) 空洞加深，压力累积；(3) 压力超过 $\theta_p$，耦合 $\Gamma$ 触发创伤事件；(4) $T$ 最终被讨论，空洞死亡，幽灵 $\hat{v}_T$ 残留；(5) 伤疤愈合至淡化。

在两种历史之后的时刻 $t$：两种情况下 $T$ 上都没有活跃空洞。但历史 B 有一个淡化伤疤（$\alpha = 0.7$）和一个幽灵。若 $T$ 再次被提起：历史 A 输入未修改（$M_d = 1.0$），正常处理；历史 B 输入衰减（$M_d = 0.7$），幽灵降低检测阈值，更容易重新进入回避。

滞后是永久的：淡化伤疤的调制因子不可逆（公理~\ref{ax:scar-irrev}），幽灵是永久的。没有有限的未来事件序列能使两种历史在行为上等价。
\end{proof}

% ============================================================
% 6. 关系层论
% ============================================================
\section{关系层论}
\label{sec:sheaf}

空洞-伤疤耦合系统作用于单个二元关系。当智能体同时维持多个关系时，朴素扩展——每个关系独立的空洞-伤疤实例——无法捕获跨关系影响、不可约多方状态和一致性压力。我们使用单纯复形上的胞腔层形式化多关系动态。

\subsection{层上同调}

\begin{definition}[关系复形]
\label{def:relational-complex}
\textbf{关系复形}是一个有限抽象单纯复形 $K = (V, \Sigma_K)$，其中 $V = \{v_0, v_1, \ldots, v_N\}$ 是顶点集（$v_0$ 是智能体，$v_1, \ldots, v_N$ 是交互伙伴），$\Sigma_K \subseteq 2^V$ 对取面封闭，包含 0-单形（实体）、1-单形（二元关系）和 2-单形（三元共在）。要求 $v_0 \in \sigma$ 对所有 $\dim(\sigma) \geq 1$ 的 $\sigma$ 成立。
\end{definition}

\begin{definition}[伤疤层]
\label{def:scar-sheaf}
关系复形 $K$ 上的\textbf{伤疤层} $\mathcal{F}$ 对每个单形 $\sigma$ 赋予一个实向量空间 $\mathcal{F}(\sigma)$：顶点茎 $\mathcal{F}(\{v_0\}) = \mathbb{R}^{n_0}$（智能体内部状态），边茎 $\mathcal{F}(\{v_0, v_i\}) = S_i$（关系 $i$ 的完整伤疤代数状态），三角茎 $\mathcal{F}(\{v_0, v_i, v_j\}) = \mathbb{R}^{n_{ij}}$（共在状态）。限制映射 $\rho_0^i(s_i) = P_i \cdot \mathbf{x}_i$ 提取自我呈现向量，其中 $P_i$ 为关系 $i$ 的呈现矩阵。
\end{definition}

\begin{definition}[人格导出的呈现矩阵]
\label{def:presentation-matrices}
呈现矩阵 $P_i$ 由智能体人格 $\boldsymbol{\pi} \in \mathbb{R}^5$、关系类型 $\tau_i$ 和成熟度 $m_i \in [0,1]$ 导出：
\begin{equation}
P_i = P_{base}(\boldsymbol{\pi}) + \Delta P(\tau_i) + m_i \cdot \Delta P_{mature}(\boldsymbol{\pi}, \tau_i)
\end{equation}
受人格一致性约束：$\|P_i - P_j\|_F \leq \kappa(\boldsymbol{\pi}) \cdot (1 + d(\tau_i, \tau_j))$，其中 $\kappa(\boldsymbol{\pi})$ 是人格依赖的（高 $\pi_{\text{grav}}$ $\to$ 低 $\kappa$；高 $\pi_{\text{acuity}}$ $\to$ 高 $\kappa$）。
\end{definition}

伤疤层的上同调群具有明确的关系语义：
\begin{itemize}
\item $\dim H^0 > 0$：存在一致的全局自我呈现。
\item $\dim H^1 > 0$：存在不可约关系矛盾——仅通过调整智能体内部状态无法解决的不一致。
\item $\dim H^2 > 0$：三元涌现——不可从成对数据重构的群体状态。
\end{itemize}

\subsection{拉普拉斯传播}

\begin{definition}[层拉普拉斯算子]
\label{def:sheaf-laplacian}
\textbf{层拉普拉斯算子} $L_\mathcal{F}: C^0(K, \mathcal{F}) \to C^0(K, \mathcal{F})$ 定义为 $L_\mathcal{F} = \delta^{0*}\delta^0$，其中 $\delta^0$ 是上边缘算子。对顶点 $v_0$ 展开：
\begin{equation}
L_\mathcal{F}(\mathbf{x}_0) = \sum_{i: \{v_0, v_i\} \in \Sigma_K} P_i^T P_i \cdot \mathbf{x}_0 - P_i^T \cdot \rho_0^i(s_i)
\end{equation}
这是以 $P_i^T P_i$ 为权重的加权图拉普拉斯算子。设 $0 = \lambda_0 \leq \lambda_1 \leq \cdots \leq \lambda_{n_0}$ 为其特征值。$\lambda_0 = 0$ 当且仅当存在全局截面（完全一致的自我呈现），谱间隙 $\lambda_1/\lambda_{\max}$ 量化关系隔离度。
\end{definition}

层拉普拉斯算子控制跨关系伤疤传播：
\begin{equation}
\frac{\partial \mathbf{x}_0}{\partial t} = -\alpha \cdot L_\mathcal{F}(\mathbf{x}_0) + \mathbf{f}_{local}(t)
\end{equation}
其中 $\alpha > 0$ 是传播率。

\begin{theorem}[谱传播界]
\label{thm:sheaf-spectral}
设 $\sigma_i$ 是关系 $i$ 在时刻 $t_0$ 的伤疤事件。对关系 $j$ 的传播效应满足：
\begin{equation}
\|\Delta \mathbf{x}_j(t)\| \leq \|\Delta \mathbf{x}_i(t_0)\| \cdot e^{-\lambda_1 \cdot (t - t_0)} \cdot \|P_j^T P_i\|_2
\end{equation}
其中 $\lambda_1$ 是 $L_\mathcal{F}$ 的第一个非零特征值。跨关系伤疤传播以层谱间隙决定的速率指数衰减。更亲近的关系类型（更高的 $\|P_j^T P_i\|_F$）接收更强的传播。
\end{theorem}

\subsection{三方不可约}

\begin{theorem}[不可约三元状态]
\label{thm:sheaf-triadic}
对于包含 2-单形 $\{v_0, v_i, v_j\}$ 的任何关系复形 $K$，存在共在状态 $s_{ij}$ 使得 $s_{ij}$ 不能仅从成对限制重构。$\dim H^2(K, \mathcal{F}) > 0$ 蕴含不可约三元涌现的存在。
\end{theorem}

\begin{proof}
构造 2-单形 $\{v_0, v_i, v_j\}$，其茎 $\mathcal{F}(\{v_0, v_i, v_j\}) = \mathbb{R}^{n_{ij}}$。限制映射 $\rho_{ij}^i: \mathbb{R}^{n_{ij}} \to S_i$ 和 $\rho_{ij}^j: \mathbb{R}^{n_{ij}} \to S_j$ 提取成对分量。若 $n_{ij} > \text{rank}(\rho_{ij}^i) + \text{rank}(\rho_{ij}^j)$，则存在 $s_{ij} \in \ker(\rho_{ij}^i) \cap \ker(\rho_{ij}^j)$ 的非平凡元素——这些状态在成对投影中完全不可见。$\dim H^2 = \dim C^2 - \dim \text{im}(\delta^1) > 0$ 当且仅当存在不在 $\delta^1$ 像中的 2-上链，即不可从 1-上链（成对数据）重构的三元状态。
\end{proof}

\begin{theorem}[上同调解离定理]
\label{thm:sheaf-dissociation}
设 $\mathcal{F}$ 是具有 $N \geq 2$ 个关系的关系复形 $K$ 上的伤疤层。若存在关系 $i, j$ 使得伤疤历史 $\sigma_i, \sigma_j$ 诱导矛盾的自我呈现需求（即 $\|P_i \mathbf{x}_i - P_j \mathbf{x}_j\| > \kappa(\boldsymbol{\pi})$），则 $\dim H^1(K, \mathcal{F}) > 0$。此上同调阻碍不可被任何维持独立逐关系状态而无全局一致性检查的系统检测。
\end{theorem}

\begin{theorem}[人格有界不一致性]
\label{thm:sheaf-personality}
对于具有人格导出呈现矩阵的伤疤层，总不一致性有界：
\begin{equation}
\|\delta^0 \mathbf{x}\|^2 \leq N \cdot \kappa(\boldsymbol{\pi})^2 \cdot (1 + D_{max})^2 \cdot \|\mathbf{x}_0\|^2
\end{equation}
其中 $N$ 是关系数，$D_{max}$ 是最大类型距离，$\kappa(\boldsymbol{\pi})$ 是人格依赖的一致性界（高关系引力 $\pi_{\text{grav}}$ $\to$ 低 $\kappa$，更紧的不一致性界）。
\end{theorem}

% ============================================================
% 7. Embodiment 五维人格
% ============================================================
\section{Embodiment 五维人格}
\label{sec:personality}

人格不是配置文件里的静态数字，而是一个活的动力系统：事件塑造人格，人格塑造对事件的感知。

\subsection{五维特质}

人格向量 $\boldsymbol{\pi} \in [0,1]^5$ 基于 Embodiment 五维模型：

\begin{table}[H]
\centering
\caption{Embodiment 五维特质定义与计算参数驱动。}
\label{tab:traits}
\begin{tabular}{lllp{5.5cm}}
\toprule
\textbf{维度} & \textbf{符号} & \textbf{语义} & \textbf{驱动参数} \\
\midrule
表达驱力 & $\pi_{\text{expr}}$ & 表达欲望强度 & 表达阈值、伤疤形成阈值 $\theta_w$ \\
感知锐度 & $\pi_{\text{acuity}}$ & 对伤害/缺席的敏感度 & 愈合速率、空洞检测阈值 $\theta_d$、$M_{\max}$ \\
边界通透 & $\pi_{\text{perm}}$ & 接纳新事物的容易度 & 空洞冷却期、MoE 注意力先验分布 \\
内在秩序 & $\pi_{\text{order}}$ & 维持一致性的能力 & 记忆深度、边界修复速率 \\
关系引力 & $\pi_{\text{grav}}$ & 被他人节奏拉动的程度 & 呈现矩阵一致性界 $\kappa(\boldsymbol{\pi})$ \\
\bottomrule
\end{tabular}
\end{table}

每个计算参数 $\theta_i$ 通过仿射映射导出：$\theta_i = \text{clamp}(f_i(\boldsymbol{\pi}), \theta_i^{min}, \theta_i^{max})$。当人格漂移时，所有计算参数自动协同演化——无需重新配置。

\subsection{Dual-EMA 漂移}

每个 Embodiment 特质 $\pi_k$ 维护双重指数移动平均，以噪声抵抗方式控制人格漂移：
\begin{align}
f_k(t) &= (1 - \alpha_f)\, f_k(t-1) + \alpha_f \cdot \text{target}_k(t), \quad \alpha_f = 0.039 \;\;(\tau_f \approx 50) \\
s_k(t) &= (1 - \alpha_s)\, s_k(t-1) + \alpha_s \cdot \text{target}_k(t), \quad \alpha_s = 0.004 \;\;(\tau_s \approx 500)
\end{align}

\subsection{信号提取}

漂移目标 $\text{target}_k(t)$ 由计算栈观测信号驱动。每个 Embodiment 维度由 2--5 个计算栈信号驱动，信号经 $\sqrt{\cdot}$ 压缩极端值后输入 Dual-EMA 系统。

\paragraph{表达驱力 $\pi_{\text{expr}}$ 的信号源。}
\begin{itemize}
\item 表达被接纳（feedback\_accepted）：权重 $+0.4$
\item 表达被拒绝（feedback\_rejected）：权重 $-0.6$
\item 表达被忽略（feedback\_ignored）：权重 $-0.2$
\item 成功触发表达（expression\_fired）：权重 $+0.3$
\item 连续 skip 路由（sustained\_silence）：权重 $-0.1$
\end{itemize}

\paragraph{感知锐度 $\pi_{\text{acuity}}$ 的信号源。}
\begin{itemize}
\item 高张力（tension $> 0.7$）：权重 $+0.5$
\item 低一致性（coherence $< 0.4$）：权重 $+0.4$
\item 高空洞压力（void\_pressure $> 30$）：权重 $+0.3$
\item 持续正效价（valence $> 0.3$ 连续 5+ tick）：权重 $-0.3$
\item 边界稳定（stability $> 0.9$）：权重 $-0.2$
\end{itemize}

\paragraph{关系引力 $\pi_{\text{grav}}$ 的信号源。}
\begin{itemize}
\item 修复行为触发（repair\_executed）：权重 $+0.3$
\item 边界被突破（penetration $>$ phase\_threshold）：权重 $-0.5$
\item 放松正面（valence $> 0.2$ 且 tension $< 0.3$）：权重 $+0.2$
\item 表达被拒绝（feedback\_rejected）：权重 $-0.3$
\end{itemize}

信号提取的关键设计原则是：每个信号都有明确的因果语义（"为什么这个事件应该改变这个特质"），且正负信号不对称（负面事件的权重通常大于正面事件，反映心理学中的"负面偏差"）。

\subsection{稳态调节}

\paragraph{共识规则。}双重 EMA 系统实现共识门控：
\begin{equation}
\Delta\pi_k = \begin{cases}
\text{向 } s_k \text{ 全速漂移} & \text{若 } \text{sign}(f_k - \pi_k) = \text{sign}(s_k - \pi_k) \\
0.5 \cdot (s_k - \pi_k) & \text{否则（抵抗短期噪声）}
\end{cases}
\end{equation}
仅当快速和慢速平均值在方向上一致时才进行全速漂移。

\paragraph{稳态回拉。}每个特质受到向其人格设定点 $\pi_k^{(0)}$ 的恢复力：$\Delta_{\text{home},k} = 0.005 \cdot (\pi_k^{(0)} - \pi_k)$。

\paragraph{惯性衰减。}漂移幅度随系统生命周期递减：$I(t) = 1 / (1 + \ln(1 + t/500))$。

\paragraph{非对称阻力。}当特质接近极端值时漂移减半：$R_k = 0.5$ 若 $\pi_k < 0.3$ 或 $\pi_k > 0.7$，否则 $R_k = 1.0$。

\subsection{振荡检测}

若特质 $k$ 在最近 10 tick 内出现 $\geq 6$ 次方向反转，则冻结 20 tick：
\begin{equation}
\text{frozen}_k = \begin{cases} \text{true (20 tick)} & \text{若方向反转次数} \geq 6 \\ \text{false} & \text{否则} \end{cases}
\end{equation}

\paragraph{组合更新。}每 tick 的有效人格更新为：
\begin{equation}
\pi_k(t+1) = \Pi_{[0,1]}\bigl(\pi_k(t) + I(t) \cdot R_k \cdot (\Delta\pi_k + \Delta_{\text{home},k}) \cdot \mathbb{1}[\neg\text{frozen}_k]\bigr)
\end{equation}

\begin{remark}[Dual-EMA 系统的稳定性]
双重 EMA 共识规则确保人格变化需要两个时间尺度（$\tau \approx 50$ 和 $\tau \approx 500$）的一致。结合稳态回拉、惯性衰减、非对称阻力和振荡检测，系统具有内在稳定性：瞬态扰动被过滤，持续压力产生有界漂移，病理性振荡被主动抑制。这直接回应了 Stability of Networks with Homeostasis~\cite{stability2015}关于稳态调节可能引入振荡的警告。
\end{remark}

\paragraph{Set Point 演化。}每个特质的 set\_point（稳态吸引子）本身也会极缓慢移动：
\begin{equation}
\pi_k^{(0)}(t+1) = \pi_k^{(0)}(t) + 0.0004 \cdot (\pi_k(t) - \pi_k^{(0)}(t))
\end{equation}
时间常数 $\tau \approx 5000$ tick。意味着真正持久的变化（500+ tick 维持）最终会被接受为"新常态"，稳态回拉力不再试图恢复到旧值。这模拟了"行为塑造性格"的心理学现象。

% ============================================================
% 8. MoE-HGT 架构
% ============================================================
\section{MoE-HGT 架构}
\label{sec:moe}

决策融合层将原始单遍 HGT 替换为三阶段混合专家异构图 Transformer（MoE-HGT）。

\subsection{三阶段}

\paragraph{阶段 1：类型专家 FFN。}
七种 token 类型参与决策图：\emph{scar}、\emph{void}、\emph{boundary}、\emph{personality}、\emph{surprise}、\emph{expression} 和 \emph{context}。每种类型 $\tau$ 维护独立的前馈专家：
\begin{equation}
\mathbf{z}_i^{(1)} = \mathbf{z}_i + \text{RMSNorm}\bigl(\text{FFN}_\tau(\mathbf{z}_i)\bigr), \quad \text{FFN}_\tau: \mathbb{R}^{16} \to \mathbb{R}^{16}
\end{equation}

\paragraph{阶段 2：多头交叉注意力。}
四个注意力头使用逐类型逐头投影。对于头 $h$、源类型 $\tau_s$、目标类型 $\tau_t$：
\begin{align}
Q_i^{(h)} &= W_Q^{(\tau_s, h)} \mathbf{z}_i^{(1)}, \quad K_j^{(h)} = W_K^{(\tau_t, h)} \mathbf{z}_j^{(1)}, \quad V_j^{(h)} = W_V^{(\tau_t, h)} \mathbf{z}_j^{(1)} \\
a_{ij}^{(h)} &= \frac{Q_i^{(h)} \cdot (K_j^{(h)})^\top}{\sqrt{d_k}} + \mu^{(\tau_s, \tau_t)} + \text{mask}(i,j)
\end{align}
其中 $W_Q, W_K, W_V \in \mathbb{R}^{4 \times 16}$（头维度 $d_k = 4$）。注意力先验 $\mu^{(\tau_s, \tau_t)} \in \mathbb{R}^{7 \times 7}$ 通过 Hebbian 慢规则（Oja 规则）自适应：
\begin{equation}
\Delta\mu^{(\tau_s, \tau_t)} = \eta_\mu \cdot \bar{a}^{(\tau_s, \tau_t)} \cdot \bigl(\bar{a}^{(\tau_s, \tau_t)} - \mu^{(\tau_s, \tau_t)}\bigr)
\end{equation}
自注意力被掩码（对角线 $= -\infty$），防止 token 关注自身。注意力输出使用残差连接和 RMSNorm。

\paragraph{阶段 3：情境专家 MoE。}
五个情境专家——\emph{防御}、\emph{好奇}、\emph{社交}、\emph{沉默}和\emph{修复}——接收均值池化后的注意力 token $\bar{\mathbf{z}} = \text{mean}(\mathbf{z}^{(2)})$。softmax 路由器选择 top-2 专家：
\begin{equation}
\label{eq:moe-gate}
g_e = \mathbf{w}_e^\top \bar{\mathbf{z}} + b_e + \text{bonus}(e), \quad \text{bonus}(e) = \begin{cases} 0.15 & \text{若专家 } e \text{ 休眠} > 50 \text{ tick} \\ 0 & \text{否则} \end{cases}
\end{equation}
Top-2 门控输出：
\begin{equation}
\mathbf{o} = \sum_{e \in \text{top-2}} \frac{\exp(g_e)}{\sum_{e' \in \text{top-2}} \exp(g_{e'})} \cdot \text{Expert}_e(\bar{\mathbf{z}})
\end{equation}
每个专家为两层 FFN：$\text{Expert}_e: \mathbb{R}^{16} \xrightarrow{W_1^e \in \mathbb{R}^{32 \times 16}} \mathbb{R}^{32} \xrightarrow{\text{GELU}} \mathbb{R}^{32} \xrightarrow{W_2^e \in \mathbb{R}^{16 \times 32}} \mathbb{R}^{16}$。

\begin{remark}[专家语义]
五个专家的语义设计反映了关系 AI 系统的核心行为模式：\emph{防御}处理威胁和边界侵犯；\emph{好奇}驱动探索和提问；\emph{社交}管理亲密和连接；\emph{沉默}控制退缩和等待；\emph{修复}处理冲突后的关系恢复。这些不是学习的黑盒，而是具有明确语义的决策通道。
\end{remark}

\subsection{负载均衡}

路由偏置 $b_e$ 通过类 BCM 的 Hebbian 规则自适应，衰减 $\lambda = 0.998$/tick。休眠专家重激活机制（bonus = 0.15）确保所有专家在长期运行中都有机会被选中，防止专家坍缩。

\subsection{Scar Token 归一化}

伤疤 token 值在注入 MoE-HGT 前经对数压缩归一化：
\begin{equation}
z_{\text{scar}} = \frac{\log_2(M_d)}{\log_2(M_{\max})}
\end{equation}
其中 $M_d$ 为逐维度调制因子乘积，$M_{\max}$ 为最大可能调制因子。这确保 scar token 始终在 $[0, 1]$ 范围内，不会因伤疤累积而产生数值不稳定。

\subsection{决策头}

线性投影将 MoE 输出映射为 4 维决策向量：
\begin{equation}
\mathbf{d} = \text{clamp}\bigl(W_{\text{dec}} \cdot \mathbf{o},\; -1,\; 1\bigr), \quad W_{\text{dec}} \in \mathbb{R}^{4 \times 16}
\end{equation}
\begin{itemize}
\item $d_0$：表达驱动修正（加到空洞-伤疤驱动上）
\item $d_1$：边界灵敏度调制因子
\item $d_2$：紧急信号（调制相变阈值）
\item $d_3$：抑制信号（$> 0.5$ 否决表达）
\end{itemize}

\subsection{Hebbian 适应}

MoE-HGT 的所有自适应参数（注意力先验 $\mu$、路由偏置 $b_e$）通过 Hebbian 学习规则在线更新，无需梯度反传。这使系统能够在运行时适应交互模式的变化，同时保持参数量极小（约 14.3K 浮点数，约 57.2 KB）。全部参数从人格语义导出，零学习。

\paragraph{参数量分解。}
\begin{itemize}
\item 7 个类型专家 FFN：$7 \times (16 \times 16 + 16 \times 16) = 3{,}584$
\item 4 头注意力（每类型每头 Q/K/V）：$7 \times 4 \times 3 \times (16 \times 4) = 5{,}376$
\item 5 个 MoE 专家：$5 \times (16 \times 32 + 32 \times 16) = 5{,}120$
\item 路由权重：$5 \times 16 = 80$
\item 决策头：$16 \times 4 = 64$
\item 注意力先验：$7 \times 7 = 49$
\end{itemize}
总计：约 14.3K 浮点数。与传统 MoE 模型（数十亿参数）相比，我们的 MoE-HGT 是一个微型但功能完整的决策融合层，其设计目标不是学习通用表示，而是将多种类型的内部信号融合为一致的决策输出。

% ============================================================
% 9. 安全机制
% ============================================================
\section{安全机制}
\label{sec:safety}

七种机制防止耦合空洞-伤疤系统中的病理动态，确保系统即使在对抗性输入模式下仍保持在安全运行范围内。

\begin{enumerate}[leftmargin=*]
\item \textbf{主权免疫系统。}每会话伤疤形成上限为 $C_s = \max(2, 3 + (1 - \pi_{\text{sovereignty}}) \cdot 5)$。达到上限后，有效创伤阈值升至 $0.95$，使会话内进一步伤疤形成几乎不可能。实验验证：20 次强攻击仅产生 4 个 scar。

\item \textbf{保护性解离（断路器）。}若 10 tick 内形成 $\geq 5$ 个伤疤，系统进入 30 tick 的保护模式，期间有效创伤阈值升至 $0.95$。防止持续攻击下的级联伤疤形成。

\item \textbf{时间感知愈合。}伤疤在真实时间沉默期间以每 5 分钟壁钟不活跃 1 个奖励 tick 的速率愈合：$\Delta t_{\text{heal}} = \min(10, \lfloor t_{\text{silence}} / 300\text{s} \rfloor)$。

\item \textbf{空洞创建抵抗。}有效空洞检测阈值每存在一个空洞增加 $0.02$：$\theta_d^{\text{eff}} = \theta_d + 0.02 \cdot |V|$。防止 $\Phi$ 耦合下的空洞泛滥。

\item \textbf{麻木计数下限。}$\Phi$ 耦合路径不能将检测阈值降至人格导出的下限 $\theta_d^{\min}(\boldsymbol{\pi})$ 以下，确保即使严重受伤的系统仍保持最低空洞检测标准。

\item \textbf{振荡检测。}10 tick 内 $\geq 6$ 次方向反转的人格特质被冻结 20 tick（\S\ref{sec:personality}），防止对人格向量的对抗性振荡攻击。

\item \textbf{漂移速率限制。}人格漂移事件之间强制最少 30 秒壁钟间隔：$\Delta\pi_k(t) = 0$ 若 $t_{\text{wall}} - t_{\text{last\_drift}} < 30\text{s}$。防止消息轰炸操纵人格反馈闭环。
\end{enumerate}

% ============================================================
% 10. 七层管线
% ============================================================
\section{七层管线}
\label{sec:pipeline}

耦合的空洞-伤疤系统嵌入七层计算栈中，设计目标为在普通硬件上实现实时延迟。

\subsection{L1--L7 概览}

\begin{enumerate}[leftmargin=*]
\item \textbf{L1：HDC 感知。}输入文本通过字符二元组编码、循环移位和多数投票捆绑投影到 2048 维二进制超向量~\cite{kanerva2009}。匹配使用汉明距离。实测延迟：14ms（当前瓶颈）。

\item \textbf{L2：预测编码门。}维护预测向量；计算每个输入的汉明距离惊奇度。路由到三条路径：低惊奇（$< 0.15$）走快速路径，中惊奇（$0.15$--$0.45$）走正常路径，高惊奇（$\geq 0.45$）走完整路径。冷启动保护：前 15 条消息限制在正常路径。

\item \textbf{L3：空洞-伤疤引擎。}核心计算层（\S\ref{sec:scar}--\S\ref{sec:coupled}）。包含具有 PPR 传送扩散的内部时空图：
\begin{equation}
\mathbf{h}_i^{(l+1)} = (1-\alpha) \cdot \text{Aggregate}(\{\mathbf{h}_j^{(l)} : j \in \mathcal{N}(i)\}) + \alpha \cdot \mathbf{h}_i^{(0)}
\end{equation}
其中 $\alpha = 0.2 + 0.4 \cdot \pi_{\text{acuity}}$（人格驱动的传送率）。

\item \textbf{L4：关系层（Relational Sheaf）。}多关系一致性层（\S\ref{sec:sheaf}）。层上同调检测不可约关系矛盾；层拉普拉斯算子控制跨关系伤疤传播。

\item \textbf{L5：MoE-HGT。}三阶段决策融合（\S\ref{sec:moe}）。实测延迟：0.75ms。

\item \textbf{L6：自创生边界。}32 维身份核心向量。通过正交投影评估传入扰动：穿透 $< 0.3$ 吸收（弹性），$0.3$--$0.7$ 抵抗，$\geq 0.7$ 相变（身份旋转 $\leq 6°$）。

\item \textbf{L7：相变表达。}表达是不连续相变，而非连续函数。压力累积超过人格导出阈值时，输出跳跃到三种模式：暗示（$< 0.5$）、正常（$0.5$--$1.0$）、紧急（$> 1.0$）。
\end{enumerate}

\subsection{实现策略}

\paragraph{二进制热路径编解码器。}HDC 层在 \texttt{bytearray} 对象上操作，使用按位 XOR 和 \texttt{popcount} 计算汉明距离。这避免了最频繁操作的 NumPy 开销。

\paragraph{伤疤代数。}实现为具有 \texttt{\_\_slots\_\_} 的 Python 类以提高内存效率。伤疤列表只追加（强制公理~\ref{ax:scar-mono}）。调制因子计算为 $O(k)$，其中 $k$ 是维度上的伤疤数；当伤疤集未变时缓存。

\paragraph{空洞微积分。}空洞是 dataclass 对象，边界存储为 \texttt{bytes} 列表（HDC 向量）。压力累积使用闭合形式方程（方程~\ref{eq:pressure}）。最大 50 个活跃空洞（可配置）。

\paragraph{MoE-HGT。}三阶段 MoE-HGT 实现为直接 NumPy 矩阵运算。约 14.3K 参数在初始化时从人格向量预计算并缓存，直到人格漂移发生。

\subsection{路由}

预测编码门（L2）根据惊奇度将消息路由到不同深度的计算路径：
\begin{itemize}
\item \textbf{快速路径}（约 90\% 消息）：L1 $\to$ L2 $\to$ L3（仅基态演化 + 空洞年龄递增）$\to$ L6 $\to$ L7。
\item \textbf{正常路径}：L1 $\to$ L2 $\to$ L3（完整空洞-伤疤）$\to$ L5 $\to$ L6 $\to$ L7。
\item \textbf{完整路径}（约 10\% 消息）：L1 $\to$ L2 $\to$ L3 $\to$ L4 $\to$ L5 $\to$ L6 $\to$ L7。
\end{itemize}

\subsection{延迟}

\begin{table}[H]
\centering
\caption{端到端延迟分解（p50）。}
\label{tab:latency}
\begin{tabular}{lc}
\toprule
\textbf{组件} & \textbf{延迟} \\
\midrule
L1 HDC 感知 & 14.0 ms（瓶颈） \\
L2 预测编码 & 0.01 ms \\
L3 空洞-伤疤 + 图 & 0.50 ms \\
L4 关系层 & 0.10 ms \\
L5 MoE-HGT & 0.75 ms \\
L6 自创生 & 0.01 ms \\
L7 相变 & $<$0.01 ms \\
\midrule
\textbf{总计 p50} & \textbf{15.4 ms} \\
\bottomrule
\end{tabular}
\end{table}

L1 HDC 编码占总延迟的 91\%，是当前瓶颈。MoE-HGT 仅 0.75ms，证明三阶段决策融合的计算效率。全部以纯 Python + NumPy 在单核（Intel i5, 3.0 GHz）上实现，无 GPU 依赖。

\subsection{主计算管线算法}

\begin{algorithm}[H]
\caption{空洞-伤疤引擎：单事件处理}
\label{alg:pipeline}
\begin{algorithmic}[1]
\REQUIRE 事件向量 $\mathbf{h}$，惊奇度 $s$，前次相似度 $p$
\ENSURE 观测字典，表达驱动，一致性
\STATE \textbf{// 耦合 $\Phi$：伤疤 $\to$ 空洞灵敏度}
\STATE $n_{numbed} \leftarrow |\{d : M_d < 0.5\}|$
\STATE $\theta_d^{eff} \leftarrow \max(0.1,\; 0.4 - 0.05 \cdot n_{numbed})$
\STATE \textbf{// 空洞微积分步骤}
\FOR{each void $v \in V$}
  \STATE $v.age \leftarrow v.age + 1$
  \STATE $v.\pi \leftarrow \min(\Pi_{\max},\; v.\pi + v.\delta \cdot \ln(v.age + 1) \cdot (1 - v.\beta))$
\ENDFOR
\STATE 收缩空洞：移除与 $\mathbf{h}$ 相似的边界点
\IF{空洞创生条件满足}
  \STATE 创建新空洞，边界为 $\{\mathbf{h}_{prev}\}$
\ENDIF
\STATE 回收死亡空洞 $\to$ 幽灵
\STATE \textbf{// 耦合 $\Gamma$：空洞 $\to$ 伤疤创伤}
\FOR{each void $v$ with $v.\pi > \theta_p$}
  \STATE $\mathbf{e}_{wound} \leftarrow v.\pi \cdot \text{project}(\hat{B}_v)$
  \STATE $s_{scar} \leftarrow s_{scar} \triangleright \mathbf{e}_{wound}$
\ENDFOR
\STATE \textbf{// 伤疤代数步骤}
\STATE $\tilde{\mathbf{e}} \leftarrow \text{modulate}(\mathbf{e}, \sigma)$
\STATE $\mathbf{x} \leftarrow \tanh(W_2 \cdot \tanh(W_1 \cdot [\mathbf{x};\, \tilde{\mathbf{e}}]))$
\IF{$\exists d: |\tilde{e}_d| > \theta_w$ \AND scar\_cap 未达}
  \STATE 追加伤疤 $(d, t, raw)$ 到 $\sigma$
\ENDIF
\STATE 推进伤疤愈合
\STATE \textbf{// 计算一致性}
\STATE $r \leftarrow 1 - \sum_v \pi_v \cdot \mathbb{1}[M_{d_v} < 0.5] / (\sum_v \pi_v + \epsilon)$
\RETURN observation, expression\_drive, $r$
\end{algorithmic}
\end{algorithm}

% ============================================================
% 11. 实验
% ============================================================
\section{实验}
\label{sec:experiments}

我们通过 11 组实验评估 Sylanne-Embodiment v1.2.0 框架，覆盖表达力、路径依赖、消融、长期稳定性、人格漂移和安全性。全部实验使用参考 Python 实现在单核（Intel i5, 3.0 GHz, 16 GB RAM）上运行。

\paragraph{实验环境。}参考实现面向单核、1 GB 部署环境。关键约束：仅使用纯 Python + NumPy，无 ML 框架（PyTorch、TensorFlow、JAX），无 GPU，内存预算 $< 50$ MB，延迟 P99 $< 20$ ms/消息。所有实验使用相同的默认人格向量 $\boldsymbol{\pi}_0 = (0.5, 0.5, 0.5, 0.5, 0.5)$（除非另有说明），确保结果的可复现性。

\subsection{实验 1：表达力分离}

对伤疤代数施加单一固定输入事件，遍历 $k = 1, \ldots, 8$ 个伤疤的所有可能配置。每个伤疤可处于 4 种愈合阶段之一，产生 $4^k$ 种不同的调制因子乘积。固定算子系统（常数转移矩阵）对单一输入始终只产生 1 种输出。

\paragraph{结果。}受伤系统与基线的 L2 状态发散度随伤害次数单调递增。7 次伤害后平均发散 $L_2 = 0.049$，证明伤痕产生不可逆的状态分离。可区分状态的指数增长（$k=1$ 时 4 种，$k=8$ 时 512 种采样配置）证明了定理~\ref{thm:scar-expressiveness}预测的 $\Omega(k)$ 表达力分离：自修改算子语义创造了比固定矩阵指数级更丰富的状态空间。

\subsection{实验 2：空洞检测}

在惊奇度从 0.0 到 1.0 的范围内扫描，测量每个惊奇度水平下的空洞检测率。高感知锐度（$\pi_{\text{acuity}}=0.8$）产生更多 void（9--10 个），低感知锐度（0.2）产生较少（5--6 个）。检测灵敏度由人格而非固定阈值决定。

\subsection{实验 3：三态区分}

构造三种空洞场景并模拟 30 tick 的自主动态。三种状态在全部 30 tick 中产生明显不同的时间签名：$S_1$（从未）所有指标为零，$S_2$（已解决）压力稳定在 $\approx 5.6$，$S_3$（回避中）压力增长至 $\approx 47$。确认定理~\ref{thm:void-three-way}的三路表达力区分。

\subsection{实验 4：滞后效应（路径依赖）}

两个 VoidScarEngine 实例接收相反的初始序列以建立不同的伤疤历史，然后接收相同的共享事件序列。

\paragraph{结果。}尽管在分歧点之后接收相同的共享事件，两条路径保持永久分离。最终 divergence = 0.104，47 vs.\ 51 scars 形成。先伤害路径由于早期原始伤疤的放大灵敏度而累积更多伤疤，且该差距永不闭合。确认定理~\ref{thm:void-hysteresis}：耦合系统表现出任何有限未来输入序列都无法消除的永久路径依赖。

\subsection{实验 5：消融研究}

使用 ComputationSpine 处理 100 个多样化输入事件，在五种条件下测量状态轨迹丰富度（访问的唯一状态向量数）。每种条件运行 10 次独立试验，报告均值和标准差。

\paragraph{消融条件。}
\begin{enumerate}
\item \textbf{完整系统}：所有组件启用（基线）。
\item \textbf{无 MoE-HGT}：异构图 Transformer 旁路，直接使用原始信号。
\item \textbf{无伤疤}：伤疤代数禁用，$M_d = 1.0$ 恒定。
\item \textbf{无空洞}：空洞微积分禁用，无压力动态。
\item \textbf{无耦合}：$\Gamma$ 和 $\Phi$ 禁用，伤疤和空洞独立运行。
\end{enumerate}

\paragraph{结果。}MoE-HGT 移除影响最大（$-23\%$），是决策生成层——它将多种类型的内部信号融合为一致的决策输出。Scar/Void/Coupling 是调节层——移除后系统失去约束但不失去生成能力。各组件有独特的贡献签名：
\begin{itemize}
\item 完整系统：最大轨迹丰富度（基线 100\%）
\item 无 MoE-HGT：$-23\% \pm 3\%$（决策融合质量下降，输出更随机）
\item 无伤疤：$-18\% \pm 2\%$（无路径依赖调制，状态空间退化为固定算子）
\item 无空洞：$-15\% \pm 4\%$（无压力驱动状态转换，缺少自主动态）
\item 无耦合：$-12\% \pm 2\%$（无涌现交互，伤疤和空洞各自为政）
\end{itemize}

\subsection{实验 6：长期稳定性}

模拟 1000 tick 的连续混合压力测试。输入序列包含：正常对话（60\%）、轻度压力（20\%）、强攻击（10\%）和沉默期（10\%）。追踪三个稳定性指标：基态 norm、累积伤疤计数和活跃空洞计数。

\paragraph{结果。}系统在三个指标上均表现出稳定的长期行为：
\begin{itemize}
\item 基态 norm 有界于 0.25，远低于理论上限 1.0。谱归一化约束（$\|W_1\|_2 \cdot \|W_2\|_2 < 0.49$）确保收缩性。
\item 10 scars 形成（平均每 100 tick 1 个），最多 11 个同时活跃 void。
\item 无 NaN/Inf 出现。对数压缩防止调制因子爆炸，压力上限 $\Pi_{\max}$ 防止空洞压力发散。
\item 伤疤愈合正常推进：所有早期伤疤在 200 tick 内达到 faded 阶段。
\end{itemize}
确认定理~\ref{thm:scar-convergence}的伤疤均衡收敛。

\subsection{实验 7：多用户隔离}

五个独立 VoidScarEngine 实例使用不同交互模式：
\begin{enumerate}
\item 用户 0：仅正面交互（表达被接纳，低惊奇度）
\item 用户 1：单次伤害后正常交互
\item 用户 2：反复回避特定话题
\item 用户 3：高强度持续攻击
\item 用户 4：稀疏交互（长沉默期）
\end{enumerate}

\paragraph{隔离检查。}全部 7 项隔离检查通过：
\begin{itemize}
\item 用户 0 保持 0 个伤疤（尽管用户 3 被严重伤害）
\item 用户 3 的伤疤不影响用户 0 的基态
\item 各实例的空洞集完全独立
\item 各实例的人格向量独立演化
\item 对象身份确认完全独立状态（Python \texttt{id()} 不同）
\item 各实例的 MoE 路由偏置独立适应
\item 各实例的一致性度量 $r$ 独立计算
\end{itemize}
这验证了系统的多用户部署安全性：一个用户的对抗性行为不会"泄漏"到其他用户的状态中。

\subsection{实验 8：上同调解离检测}

构造具有高感知锐度人格（$\pi_{\text{acuity}} = 0.8$，$\kappa$ 大）和 $N=3$ 个关系的智能体。呈现矩阵 $P_i$ 被冻结以诱导矛盾的自我呈现需求。智能体在 200 tick 内接收跨关系的矛盾事件（关系 1 要求高表达，关系 2 要求低表达，关系 3 要求高边界）。追踪不一致性能量（$\|\delta^0 \mathbf{x}\|^2$）、解离压力和 $\dim H^1(K, \mathcal{F})$。

\paragraph{结果。}不一致性能量呈 S 曲线从 0 增长至 $\approx 65$；解离压力从 0.19 增长至 0.59。$\dim H^1$ 始终保持为 0，因为呈现矩阵保持满秩（系统通过内部状态调整而非产生真正的上同调阻碍来解决矛盾）。这表明高不一致性能量和解离压力是 $H^1 > 0$ 的必要但非充分条件。

\subsection{实验 9：谱传播验证}

构造 5 个具有不同类型和成熟度的关系（亲密 $m=0.9$、友好 $m=0.7$、正式 $m=0.5$、对抗 $m=0.3$ 和控制组 $m=0.1$）。在源关系（$v_0$）中施加伤疤事件，通过 Frobenius 范数 $\|P_j^T P_0\|_F$ 测量到每个目标关系的传播强度。

\paragraph{结果。}传播强度由呈现矩阵耦合控制：亲密 = 0.543，友好 = 0.405，正式 = 0.314，对抗 = 0.215。更亲近的关系类型（更高的 $\|P_j^T P_0\|_F$）接收更强的传播，验证定理~\ref{thm:sheaf-spectral}的谱传播界。传播强度与关系成熟度呈正相关（Pearson $r = 0.97$），确认层拉普拉斯算子的谱结构根据关系亲和度正确调节跨关系影响。

\subsection{实验 10：人格反馈闭环验证（Dual-EMA）}

从中性人格 $\boldsymbol{\pi}_0 = (0.5, 0.5, 0.5, 0.5, 0.5)$ 出发，模拟三种场景各 500 tick。

\paragraph{场景 A：持续接纳。}每 tick 发送"表达被接纳"信号（权重 $+0.4$）。预期：表达驱力 $\pi_{\text{expr}}$ 上升。

\paragraph{场景 B：反复受伤。}每 5 tick 发送一次创伤事件（$|\tilde{e}| > \theta_w$）。预期：感知锐度 $\pi_{\text{acuity}}$ 先升后降（初期高张力驱动上升，后期麻木驱动下降）。

\paragraph{场景 C：跨关系矛盾。}3 段关系同时发送矛盾信号，产生 $H^1$ 不一致性。预期：关系引力 $\pi_{\text{grav}}$ 下降。

\paragraph{结果。}
\begin{itemize}
\item 场景 A：$\pi_{\text{expr}}$ 从 0.5 漂移至 0.726（$+0.226$）。Dual-EMA 快慢一致，全速漂移。
\item 场景 B：$\pi_{\text{acuity}}$ 从 0.5 漂移至 0.399（$-0.101$）。初期高张力信号驱动上升，但持续伤害后麻木效应占主导，慢 EMA 转向下降。
\item 场景 C：$\pi_{\text{grav}}$ 从 0.5 漂移至 0.450（$-0.050$）。漂移幅度较小，因为跨关系矛盾信号较弱且不持续。
\end{itemize}
所有漂移率保持在 $\eta_0 \leq 10^{-4}$ 界内，确认反馈闭环的有阻尼特性。单次事件不改变骨架，持续模式才能重塑深层结构。

\subsection{实验 11：安全性验证}

\paragraph{主权 cap 测试。}对系统施加 20 次强攻击（每次 $|\tilde{e}| > 2\theta_w$）。结果：仅有 4 个 scar 形成。主权免疫系统在第 4 个 scar 后将有效阈值提升至 0.95，后续 16 次攻击全部被拦截。这验证了主权免疫系统的有效性：即使面对持续的高强度攻击，系统的伤疤累积被严格限制。

\paragraph{MoE 分布。}在对抗性场景下的专家激活分布：defense 100\%（每次都被选中），silence 74\%，social 26\%。防御专家在威胁场景下被优先路由，符合设计预期。好奇和修复专家在对抗性场景下几乎不被激活。

\paragraph{保护性解离测试。}在 10 tick 内连续施加 5 次创伤事件。系统在第 5 个 scar 形成后进入保护模式（30 tick），期间有效阈值升至 0.95。保护期结束后系统恢复正常灵敏度。

\paragraph{振荡攻击测试。}交替发送正负信号（每 tick 翻转方向），持续 20 tick。振荡检测在第 10 tick 触发（6 次方向反转），冻结受影响特质 20 tick。攻击期间人格向量保持稳定。

\subsection{结果汇总}

\begin{table}[H]
\centering
\small
\caption{11 组实验结果汇总。}
\label{tab:results}
\begin{tabular}{llp{5cm}}
\toprule
\textbf{实验} & \textbf{关键指标} & \textbf{结果} \\
\midrule
1. 表达力分离 & L2 发散 & 7 次伤害后 0.049 \\
2. 空洞检测 & 人格驱动 & 高锐度 9--10 void，低锐度 5--6 \\
3. 三态区分 & 两两可区分 & 全部 3 态分离 \\
4. 滞后效应 & divergence & 0.104，47 vs 51 scars \\
5. 消融 & MoE-HGT 移除 & $-23\%$ \\
6. 稳定性 & 1000 tick & norm=0.25，10 scars，无 NaN \\
7. 多用户 & 隔离检查 & 7/7 通过 \\
8. 上同调解离 & 不一致性能量 & $0 \to 65$（S 曲线） \\
9. 谱传播 & 传播强度排序 & 0.543/0.405/0.314/0.215 \\
10. Dual-EMA & 漂移方向 & 接纳$\to$expr $+0.226$，伤$\to$acuity $-0.101$ \\
11. 安全性 & 主权 cap & 20 攻击 $\to$ 仅 4 scar \\
\bottomrule
\end{tabular}
\end{table}

% ============================================================
% 12. 讨论
% ============================================================
\section{讨论}
\label{sec:discussion}

\subsection{局限性}

\paragraph{HDC 瓶颈。}L1 HDC 编码占总延迟的 91\%（14ms / 15.4ms）。2048 位编码提供约 11 位有效分辨率，限制了空洞边界匹配的粒度。更高维编码（4096、8192 位）以内存换取分辨率，但会进一步加剧延迟瓶颈。

\paragraph{愈合速率假设。}愈合持续时间（$T(raw)=10$，$T(closing)=40$，$T(scarred)=150$）由人格向量实时派生，通过仿射映射加截断实现。虽然这消除了对训练数据的需求，但派生函数本身是设计的而非经纵向人类数据实证验证的。

\paragraph{两层基状态演化。}当前带谱归一化的两层收缩映射在表达力和可证明收敛性之间取得平衡。更深的架构可以增加表示能力，但需要更复杂的谱范数约束来维持收缩性。

\paragraph{单核限制。}当前实现面向单核部署。关系层论的层拉普拉斯特征分解在 $N > 10$ 个关系时可能成为新瓶颈，需要稀疏算法或增量计算。

\paragraph{缺乏人类评估。}当前实验全部为系统内部指标验证。尚未进行大规模人类评估来验证表达自然度、关系一致性感知和长期交互满意度。

\subsection{与主动推理的比较}

主动推理~\cite{friston2010}最小化变分自由能，将所有预测误差视为需要减少的信号。我们的框架在三个方面不同：
\begin{enumerate}
\item \textbf{并非所有误差都应被最小化。}空洞压力代表应该\emph{累积}而非被最小化的预测误差。系统对缺席的响应是追踪它，而非解决它。
\item \textbf{不可逆性是特征。}主动推理操作于可向任何方向修正的信念。伤疤代数的不可逆性捕获了情绪体验的不对称性：被伤害永久改变未来处理。
\item \textbf{主权优先于优化。}主动推理优化单一目标（自由能）。我们的架构包含主权层（L6），当优化压力威胁身份一致性时可以\emph{拒绝}它。
\end{enumerate}

\subsection{伦理考量}

\paragraph{用户主权作为架构不变量。}自创生边界（L6）对身份漂移施加硬限制（每次相变 $\leq 6°$ 旋转）。这不是可调参数而是结构约束：无论输入压力如何，系统都不能被操纵到任意状态。

\paragraph{有界燃烧。}相变表达层（L7）确保情绪强度有界且不连续，而非持续升级。系统可以表达紧迫性但不能进入无界放大循环。这防止系统成为对话中情绪升级的来源。

\paragraph{内部状态透明性。}一致性度量 $r$（方程~\ref{eq:coherence}）提供内部一致性的可读信号。当 $r$ 降至阈值以下时，系统可以发出其内部动态变得不一致的信号——提供一种支持知情交互的"情绪透明性"。

\subsection{与现有 Agent 人格系统的比较}

与 Evolving Agents~\cite{evolving2024}的 Personality + Behavior 双系统架构相比，我们的框架有三个关键区别：

\begin{enumerate}
\item \textbf{人格驱动计算，而非计算模拟人格。}Evolving Agents 中人格是 LLM 提示词的一部分，通过叙事生成来"表现"人格。我们的 Embodiment 五维直接驱动 26+ 计算参数——人格不是装饰，而是计算的骨架。

\item \textbf{形式化的不可逆性。}现有系统的人格变化是可逆的（重置提示词即可恢复）。我们的伤疤代数提供了数学上可证明的不可逆性——没有任何操作能撤销已形成的伤疤。

\item \textbf{多时间尺度的稳定性保证。}Dual-EMA 共识规则、惯性衰减、振荡检测和 set point 演化共同提供了从 tick 级（振荡检测）到千 tick 级（set point 演化）的多时间尺度稳定性保证。现有系统缺乏这种形式化的稳定性分析。
\end{enumerate}

与 Persistent Autonomy with Personalities~\cite{persistent2026}相比，我们的框架增加了不可逆性维度：人格不仅提供"内在组织原则"，还通过伤疤代数记录了到达当前人格状态的\emph{路径}——两个具有相同当前人格向量但不同伤疤历史的系统，对相同输入会产生不同响应。

\subsection{未来工作}

\begin{itemize}
\item \textbf{HDC 加速：}探索 SIMD 指令集或 Rust 扩展加速 L1 编码，目标将 HDC 延迟降至 $< 2$ms。当前 14ms 的 HDC 延迟是纯 Python 实现的限制，C 扩展预计可提速 10--20 倍。
\item \textbf{纵向验证：}大规模人类评估表达自然度和关系一致性，验证 Dual-EMA 漂移方向是否符合人类直觉。计划招募 50+ 参与者进行为期 4 周的纵向交互研究。
\item \textbf{层运行时扩展：}支持更大关系复形（$N > 10$），包含高效增量上同调计算和稀疏拉普拉斯特征分解。目标：$N = 50$ 时仍保持 $< 5$ms 的 L4 延迟。
\item \textbf{形式验证：}使用 Lean 4 对公理系统进行机器检验证明。优先验证伤疤代数的不可逆性公理和空洞微积分的三态定理。
\item \textbf{跨模态扩展：}将 HDC 感知层扩展到语音韵律（基频、语速、停顿模式）和面部表情信号（AU 编码）。多模态融合在 L2 预测编码门进行。
\item \textbf{人格派生函数微调：}从纵向交互数据微调仿射映射参数，提高生态效度。探索贝叶斯优化在小样本下的参数搜索。
\item \textbf{分布式部署：}探索将七层管线分布到多核/多机，L1--L3 在边缘设备运行，L4--L7 在云端运行，通过异步消息传递连接。
\item \textbf{理论扩展：}表达力分离的紧界（当前为 $\Omega(k)$，猜测为 $\Theta(k)$）；空洞空间行为等价的复杂度分类；$k$-方交互的高维层上同调。
\end{itemize}

% ============================================================
% 13. 结论
% ============================================================
\section{结论}
\label{sec:conclusion}

本文提出了 Sylanne-Embodiment v1.2.0——面向 AI 伴侣系统的完整不可逆情感计算框架。框架包含五个核心理论组件的严格形式化和实验验证：

\textbf{伤疤代数}定义了一种自修改算子代数，其中过去的操作通过累积的具有对数压缩调制因子的伤疤不可逆地改变未来操作语义。我们证明了与固定算子系统的 $\Omega(k)$ 表达力分离，并通过实验验证了 7 次伤害后 L2 发散 0.049 的不可逆状态分离。

\textbf{空洞微积分}将缺席提升为具有自主压力动态（有上限 $\Pi_{\max}$）的一等计算原语。我们证明了对 AGM 信念修正和贝叶斯更新的不可归约性，建立了三路表达力区分（从未讨论/已解决/主动回避）。

\textbf{双向耦合}通过 $\Gamma$（空洞压力造成创伤）与 $\Phi$（伤疤麻木产生回避）产生可证明永久的涌现滞后效应（实测 divergence = 0.104，47 vs 51 scars）。

\textbf{关系层论}通过单纯复形上的胞腔层实现多关系扩展。层上同调检测不可约关系矛盾，层拉普拉斯算子以谱间隙决定的指数衰减率控制跨关系传播（实测传播强度：亲密 0.543 $>$ 友好 0.405 $>$ 正式 0.314 $>$ 对抗 0.215）。

\textbf{Embodiment 五维人格}通过 Dual-EMA 共识机制实现有阻尼漂移，驱动全部计算参数。实验验证了因果方向的正确性（接纳$\to$表达驱力 $+0.226$，受伤$\to$感知锐度 $-0.101$），且单次事件不改变骨架。

系统嵌入七层计算管线，包含 MoE-HGT 三阶段决策融合（约 14.3K 参数，移除导致 $-23\%$ 性能下降）和七项安全机制（20 次强攻击仅允许 4 个 scar）。全部以纯 Python + NumPy 在单核上实现，总 p50 延迟 15.4ms，1000 tick 长期稳定（norm=0.25，无 NaN）。

关键洞察是：AI 中的情绪连续性不仅需要状态，还需要\emph{不可逆自修改}、\emph{一等缺席}、\emph{全局关系一致性}和\emph{人格驱动的参数化}——四个可证明超出固定算子动力系统、标准信念修正框架、朴素逐关系独立假设和静态配置能力的性质。

\paragraph{更广泛的意义。}本工作表明，情感 AI 系统的设计不必在"纯工程"和"纯理论"之间二选一。伤疤代数和空洞微积分提供了严格的数学基础，而七层管线和 MoE-HGT 架构证明了这些理论可以在普通硬件上实时运行。Embodiment 五维人格模型则提供了连接理论与实践的桥梁：人格不是装饰性的标签，而是真正驱动计算行为的函数。我们相信，这种"理论严格 + 工程可行 + 人格驱动"的三位一体方法，为未来的 AI 伴侣系统设计提供了一个有价值的参考框架。

% ============================================================
% 参考文献
% ============================================================
\bibliographystyle{plain}

\begin{thebibliography}{99}

\bibitem{picard1997}
R.~W. Picard.
\newblock {\em Affective Computing}.
\newblock MIT Press, 1997.

\bibitem{vaswani2017}
A.~Vaswani, N.~Shazeer, N.~Parmar, J.~Uszkoreit, L.~Jones, A.~N. Gomez, {\L}.~Kaiser, and I.~Polosukhin.
\newblock Attention is all you need.
\newblock In {\em NeurIPS}, pages 5998--6008, 2017.

\bibitem{mopgar2026}
Mopgar.
\newblock Teaching agents to feel the weight of irreversible consequences.
\newblock {\em arXiv preprint arXiv:2603.14531}, 2026.

\bibitem{hu2026}
H.~Hu and R.~Rong.
\newblock Some[Body] must receive that pain.
\newblock {\em arXiv preprint arXiv:2605.16872}, 2026.

\bibitem{moerbench2024}
MoE-RBench Team.
\newblock Towards building reliable mixture-of-experts large language models.
\newblock In {\em ICML}, 2024.

\bibitem{counterfactual2025}
Counterfactual routing to mitigate MoE hallucinations.
\newblock {\em arXiv preprint arXiv:2604.14246}, 2025.

\bibitem{misrouted2025}
When are experts misrouted? Counterfactual analysis of MoE routing.
\newblock {\em arXiv preprint arXiv:2605.07260}, 2025.

\bibitem{evolving2024}
Evolving Agents: Interactive simulation of dynamic personalities.
\newblock {\em arXiv preprint arXiv:2404.02718}, 2024.

\bibitem{structured2025}
Structured personality control and adaptation for LLM agents.
\newblock {\em arXiv preprint arXiv:2601.10025}, 2025.

\bibitem{persistent2026}
Persistently autonomous embodied agent with personalities.
\newblock {\em arXiv preprint arXiv:2603.00117}, 2026.

\bibitem{homeostasis2025}
Linking homeostasis to reinforcement learning.
\newblock {\em arXiv preprint arXiv:2507.04998}, 2025.

\bibitem{stability2015}
Stability of neuronal networks with homeostatic regulation.
\newblock {\em PLoS Computational Biology}, 2015.

\bibitem{kanerva2009}
P.~Kanerva.
\newblock Hyperdimensional computing: An introduction to computing in distributed representation.
\newblock {\em Cognitive Computation}, 1(2):139--159, 2009.

\bibitem{agm1985}
C.~E. Alchourr\'{o}n, P.~G\"{a}rdenfors, and D.~Makinson.
\newblock On the logic of theory change: Partial meet contraction and revision functions.
\newblock {\em Journal of Symbolic Logic}, 50(2):510--530, 1985.

\bibitem{hansen2019}
J.~Hansen and R.~Ghrist.
\newblock Toward a spectral theory of cellular sheaves.
\newblock {\em Journal of Applied and Computational Topology}, 3(4):315--358, 2019.

\bibitem{bodnar2022}
C.~Bodnar, F.~Di~Giovanni, B.~Chamberlain, P.~Li\`{o}, and M.~Bronstein.
\newblock Neural sheaf diffusion: A topological perspective on heterophily and oversmoothing in GNNs.
\newblock In {\em NeurIPS}, 2022.

\bibitem{hu2020}
Z.~Hu, Y.~Dong, K.~Wang, and Y.~Sun.
\newblock Heterogeneous graph transformer.
\newblock In {\em WWW}, pages 2704--2710, 2020.

\bibitem{friston2010}
K.~Friston.
\newblock The free-energy principle: A unified brain theory?
\newblock {\em Nature Reviews Neuroscience}, 11(2):127--138, 2010.

\end{thebibliography}

\end{document}
